\section{Ruido, adquisición y algoritmo de estimación}
\subsection{Modelo de ruido en unidades de potencia}
La medida simulada es $y_\ell=P_\ell(\vect p)+\varepsilon_\ell$, con ruido Gaussian independiente por piloto y desviación
\begin{equation}\label{eq:noise}
\sigma_P^2(P)=\frac{B}{\mathcal R^2}\left[
\frac{4k_BT}{R_f}+i_a^2+2q_e(\mathcal RP+I_{bg})+
\mathrm{RIN}(\mathcal RP)^2\right]+(\epsilon P)^2.
\end{equation}
La responsividad $\mathcal R$ convierte potencia en corriente; los términos entre corchetes son densidades espectrales de corriente, que se integran sobre el ancho de banda $B$ y se vuelven a expresar en W$^2$. Se incluyen ruido térmico, amplificador, shot de señal y fondo, RIN y una fluctuación multiplicativa por piloto.

\begin{table}[H]\centering\begin{tabular}{ll}\toprule
Parámetro & Valor\\\midrule
$\mathcal R$ & 0.7 A/W\\
$T$, $R_f$ & 300 K, 10 k$\Omega$\\
$i_a$ & 2 pA/$\sqrt{\mathrm{Hz}}$\\
$I_{bg}$, con media sustraída & 10 $\mu$A\\
RIN & $-155$ dB/Hz\\
$\epsilon$ & 0.5\%, independiente entre pilotos\\
$B_0$ de referencia & 100 Hz\\
Desviación sin señal, $B=B_0$ & aproximadamente 42.53 pW\\\bottomrule
\end{tabular}\caption{Hipótesis de electrónica para el PD único. No se utiliza el ancho de banda de 2 GHz del receptor de comunicaciones del artículo.}\end{table}

La componente multiplicativa no representa un error fijo de calibración: es una fluctuación independiente entre adquisiciones. Las incertidumbres persistentes por VCSEL, focal o ganancia común se tratan como perturbaciones separadas. No deben confundirse porque promediar puede reducir la primera, pero no corregir las segundas.

\subsection{Comparación justa del número de estados}
Se usa encendido secuencial, de modo que el shot de señal corresponde al emisor activo. Para una integración rectangular de duración $T_{int}$, la banda equivalente unilateral es $B=1/(2T_{int})$. Se fija
\begin{equation}
B_K=100K\ \mathrm{Hz},\qquad T_{int,K}=\frac1{200K}\ \mathrm{s}.
\end{equation}
El tiempo de integración de $25K$ pilotos es 125 ms por tile; para $225K$, 1.125 s. Así, aumentar $K$ no recibe gratuitamente más tiempo de integración.

Esta igualdad no incluye asentamiento de la lente ni procesamiento. Tampoco hace equivalente un sistema de pilotos simultáneos: encender varios emisores cambia shot, RIN e interferencia. La posible reducción de latencia mediante pilotos ortogonales requiere un modelo de adquisición distinto. Las simulaciones presuponen posición constante durante todo el barrido.

Se incluye un control especialmente importante: repetir exactamente $R=15$ mm dos y ocho veces. Ese control puede promediar la fluctuación multiplicativa independiente, pero no aporta nueva geometría. Compararlo con estados distintos permite separar promediado y diversidad real.

\subsection{Función objetivo sin información de la posición verdadera}
El generador conoce la verdad para crear datos; el estimador no la recibe. Sus pesos se calculan con las medidas:
\begin{equation}
\widehat\sigma_\ell=\sigma_P(\max(y_\ell,0)),\qquad
\widehat{\vect p}=\arg\min_{\vect p\in\mathcal B}
\sum_\ell\frac{(P_\ell(\vect p)-y_\ell)^2}{\widehat\sigma_\ell^2}.
\end{equation}
Esto es mínimos cuadrados ponderados factibles, no máxima verosimilitud exacta del modelo heteroscedástico. Las medidas negativas después de sustraer fondo se conservan en el residuo. Solo se limitan a cero al evaluar la varianza y la firma comprimida de inicialización.

Si la ganancia es desconocida, el estado es $[x,y,z,\gamma]^T$, con $g=e^\gamma$, y el residuo es $(gP-y)/\widehat\sigma$. Se usan límites $0.2\le g\le5$. Para altura conocida se omite $z$ del estado. La derivada adicional respecto de $\gamma$ es $gP$, y la de posición es $gJ$.

\subsection{Biblioteca, varios arranques y ajuste local}
El procedimiento implementado es:
\begin{enumerate}
\item Construir una biblioteca de potencias del modelo en una malla de la caja de búsqueda.
\item Ordenar los candidatos por coste ponderado. Si la ganancia es libre, perfilar su valor inicial mediante $g=\langle P,y\rangle_w/\langle P,P\rangle_w$, limitado al intervalo admisible.
\item Ordenar también por distancia entre firmas $\operatorname{asinh}(P/\sigma_{oscuro})$. Esta compresión evita que un ligero desajuste de malla de un canal fuerte descarte la región correcta.
\item Intercalar ambos órdenes y seleccionar seis semillas separadas al menos 12 cm.
\item Ejecutar \archivo{lsqnonlin}, con Jacobiano analítico, restricciones de caja y algoritmo trust-region-reflective.
\item Conservar la solución de menor coste y registrar mínimos alternativos separados más de 10 cm.
\end{enumerate}
La compresión \emph{no} cambia el modelo de observación ni el ajuste final: únicamente elige candidatos. La malla de inicialización es distinta de las posiciones aleatorias y de la malla de evaluación. Para el estudio se utilizan $35\times35$ candidatos laterales por tile, $51\times51$ para nueve tiles y nueve alturas en 3D. Las variables internas del optimizador y las semillas quedan accesibles en el struct devuelto.

No se promete encontrar el mínimo global de una función no convexa. La convergencia de \archivo{lsqnonlin} se informa por separado del error real en simulación. En un punto oscuro, un algoritmo puede converger y seguir siendo incapaz de proporcionar una localización útil.

\subsection{Regresión que comprueba la búsqueda}
En Python se encontró una observación del AP completo cuya posición verdadera era aproximadamente $(-0.1024,-1.0424,0)$ m. La búsqueda basada solo en coste ponderado arrancaba en una región errónea y terminaba a 1.396 m de la verdad, pese a una información local favorable. Se introdujo la búsqueda combinada, sin suministrar la verdad al algoritmo.

La observación concreta se conserva como referencia y se vuelve a ajustar en MATLAB. \archivo{testOriginalInitializerRegression} exige error inferior a 1 mm y coste inferior a 600. Se verifica así el mecanismo que resolvió el fallo, no solo que MATLAB ejecuta la sintaxis.

\textbf{Archivos del capítulo:} \archivo{configuracion/gob_noise.m}, \archivo{modelo/gob_sigma.m}, \archivo{estimacion/gob_estimator.m}, \archivo{estimacion/gob_candidate_costs.m}, \archivo{estimacion/gob_residual.m}, \archivo{estimacion/gob_fit.m} y \archivo{demo_localizacion.m}.
